% ===================== Chủ đề 1 -- Đại số và Phương trình -- phần c (VI): Bài 31--40 =====================
\providecommand{\aeaBox}[4]{\begin{tikzpicture}[baseline=-0.7ex,font=\small]
  \draw[Navy,line width=0.8pt,fill=tmBlue!6] (-0.66,-0.32) rectangle (0.66,0.32);
  \draw[Navy,line width=0.8pt] (0,-0.32)--(0,0.32) (-0.66,0)--(0.66,0);
  \node[inner sep=0pt] at (-0.33,0.16){$#1$};\node[inner sep=0pt] at (0.33,0.16){$#2$};
  \node[inner sep=0pt] at (-0.33,-0.16){$#3$};\node[inner sep=0pt] at (0.33,-0.16){$#4$};\end{tikzpicture}}

\clearpage
\dividerbar{Mở rộng -- Nâng cao}

% ---------------------------------------------------------------- Bài 31
\begin{probbox}{31}{\relax}
Quãng đường từ nhà bác Ba đến thành phố dài $180$ km. Bác đi xe máy với vận tốc không đổi và sau $3$ giờ thì đi được
$\dfrac{2}{3}$ quãng đường.
\begin{enumerate}[label=\alph*),leftmargin=1.6em,itemsep=1pt,topsep=2pt]
\item Tính vận tốc của xe máy.
\item Nếu giữ nguyên vận tốc đó, bác cần đi thêm bao lâu nữa thì đến thành phố?
\item Nếu bác muốn đến sớm hơn dự kiến ở câu b) $30$ phút thì trên phần đường còn lại bác phải đi với vận tốc bao nhiêu?
\end{enumerate}
\end{probbox}
\probsol{
\textbf{Phân tích đề.} \emph{Đã biết:} cả quãng đường $180$ km; trong $3$ giờ đi được $\dfrac{2}{3}$ quãng đường.
\emph{Cần tìm:} vận tốc; thời gian còn lại; vận tốc mới nếu muốn bớt $30$ phút.
\begin{recallbox}$s=v\cdot t$;\quad $v=\dfrac{s}{t}$;\quad $t=\dfrac{s}{v}$. \ \ $\dfrac{m}{n}$ của $A$ bằng $A:n\cdot m$.\end{recallbox}
\step{1}{Quãng đường đã đi và quãng đường còn lại}
Đã đi: $\dfrac{2}{3}$ của $180$ km là $180:3\cdot 2=60\cdot 2=120$ (km). Còn lại: $180-120=60$ (km).
\begin{barfigC}
\blk{0}{2}{0}{0.55}{tmBlue}{60 km}
\blk{2}{4}{0}{0.55}{tmBlue}{60 km}
\blk{4}{6}{0}{0.55}{tmOrange}{60 km}
\brcup{0}{4}{0.55}{đã đi (3 giờ)}
\brcdn{4}{6}{-0.55}{còn lại}
\end{barfigC}
\step{2}{Câu a: vận tốc}
$v=120:3=40$ (km/h).
\step{3}{Câu b: thời gian còn lại}
$t=60:40=1{,}5$ (giờ) $=1$ giờ $30$ phút.
\begin{tipbox}Nhìn sơ đồ: $3$ giờ đi được $2$ đoạn bằng nhau, vậy mỗi đoạn mất $1{,}5$ giờ. Phần còn lại đúng $1$ đoạn nên mất $1{,}5$ giờ. Không cần tính vận tốc vẫn ra!\end{tipbox}
\step{4}{Câu c: vận tốc mới}
Thời gian mới: $1$ giờ $30$ phút $-\,30$ phút $=1$ giờ. Để đi $60$ km trong $1$ giờ, vận tốc phải là $60:1=60$ (km/h).
\begin{warnbox}Không được cộng thêm vào vận tốc theo kiểu ``bớt $30$ phút thì tăng $30$ km/h''. Phải tính lại: vận tốc $=$ quãng đường $:$ thời gian mới.\end{warnbox}
\verify{$40\cdot 3=120$ km; $40\cdot 1{,}5=60$ km; $60\cdot 1=60$ km. Tổng $120+60=180$ km.}
\ans{a) $40$ km/h;\ \ b) $1$ giờ $30$ phút;\ \ c) $60$ km/h.}
}

% ---------------------------------------------------------------- Bài 32
\begin{probbox}{32}{\relax}
Một tam giác cân có cạnh đáy $10$ cm và hai cạnh bên, mỗi cạnh dài $2x$ cm. Một hình chữ nhật có chiều dài $(x+7)$ cm và chiều rộng $5$ cm.
Biết chu vi tam giác bằng chu vi hình chữ nhật.
\begin{figC}
\draw[tmBlue,line width=1.1pt,fill=tmBlue!10] (0,0)--(2.4,0)--(1.2,2.6)--cycle;
\node[below] at (1.2,0){$10$ cm};
\node[left] at (0.5,1.35){$2x$};\node[right] at (1.9,1.35){$2x$};
\draw[tmBlue,line width=0.9pt] (0.48,1.4)--(0.72,1.2); \draw[tmBlue,line width=0.9pt] (1.68,1.2)--(1.92,1.4);
\draw[tmOrange!85!black,line width=1.1pt,fill=tmOrange!12] (3.6,0) rectangle (7.2,1.3);
\node[below] at (5.4,0){$x+7$}; \node[right] at (7.2,0.65){$5$ cm};
\end{figC}
\begin{enumerate}[label=\alph*),leftmargin=1.6em,itemsep=1pt,topsep=2pt]
\item Viết biểu thức tính chu vi của mỗi hình theo $x$.
\item Tìm $x$, rồi tính độ dài các cạnh của tam giác và diện tích hình chữ nhật.
\end{enumerate}
\end{probbox}
\probsol{
\textbf{Phân tích đề.} Tam giác có ba cạnh $2x$, $2x$, $10$ (cm). Hình chữ nhật có kích thước $x+7$ và $5$ (cm). Hai chu vi bằng nhau.
\begin{recallbox}Chu vi tam giác $=$ tổng ba cạnh. Chu vi hình chữ nhật $=2\times(\text{dài}+\text{rộng})$. Diện tích hình chữ nhật $=$ dài $\times$ rộng.\end{recallbox}
\step{1}{Câu a: biểu thức chu vi}
Tam giác: $2x+2x+10=4x+10$ (cm).\quad Hình chữ nhật: $2\cdot\big((x+7)+5\big)=2(x+12)=2x+24$ (cm).
\begin{warnbox}Khi nhân $2$ vào ngoặc phải nhân cả hai số hạng: $2(x+12)=2x+24$, không phải $2x+12$.\end{warnbox}
\step{2}{Câu b: lập phương trình}
$4x+10=2x+24$.
\step{3}{Giải phương trình}
Trừ $2x$ ở cả hai vế: $2x+10=24$. Trừ $10$ ở cả hai vế: $2x=14$. Chia cho $2$: $x=7$.
\step{4}{Tính các đại lượng được hỏi}
Cạnh bên tam giác: $2x=14$ cm, nên tam giác có các cạnh $14$ cm, $14$ cm, $10$ cm.\\
Hình chữ nhật: dài $x+7=14$ cm, rộng $5$ cm, diện tích $14\cdot 5=70$ (cm$^2$).
\verify{Chu vi tam giác $14+14+10=38$ cm; chu vi hình chữ nhật $2(14+5)=38$ cm. Bằng nhau.}
\ans{a) $4x+10$ và $2x+24$;\ \ b) $x=7$; tam giác $14$ cm, $14$ cm, $10$ cm; diện tích hình chữ nhật $70$ cm$^2$.}
}

% ---------------------------------------------------------------- Bài 33
\begin{probbox}{33}{\stars{1}}
Ta dùng lại quy tắc của Bài 14:\quad \aeaBox{a}{b}{c}{d} $=a\cdot d+b\cdot c$.
\begin{enumerate}[label=\alph*),leftmargin=1.6em,itemsep=2pt,topsep=2pt]
\item Tính \aeaBox{-3}{4}{5}{2}.
\item Tìm $x$ biết \aeaBox{x}{2}{3}{4} $=$ \aeaBox{1}{x}{5}{2}.
\item Tìm tất cả các số nguyên $x$ sao cho \aeaBox{x}{x}{x}{x} $=50$.
\end{enumerate}
\end{probbox}
\probsol{
\textbf{Phân tích đề.} Quy tắc: nhân hai số trên mỗi đường chéo rồi cộng hai tích. Câu a là tính giá trị (có số âm), câu b và c là giải phương trình.
\step{1}{Câu a}
\aeaBox{-3}{4}{5}{2} $=(-3)\cdot 2+4\cdot 5=-6+20=14$.
\begin{recallbox}Tích của một số âm và một số dương là số âm: $(-3)\cdot 2=-6$. Cộng hai số khác dấu: $-6+20=20-6=14$.\end{recallbox}
\step{2}{Câu b: viết hai vế theo quy tắc}
Vế trái: $x\cdot 4+2\cdot 3=4x+6$.\quad Vế phải: $1\cdot 2+x\cdot 5=2+5x$.\quad Phương trình: $4x+6=5x+2$.
\step{3}{Giải câu b}
Trừ $4x$ ở hai vế: $6=x+2$. Trừ $2$ ở hai vế: $x=4$.
\verify{Vế trái $4\cdot 4+6=22$; vế phải $2+5\cdot 4=22$. Bằng nhau.}
\step{4}{Câu c}
\aeaBox{x}{x}{x}{x} $=x\cdot x+x\cdot x=2x^2$. Phương trình $2x^2=50$, tức là $x^2=25$.
Số nào bình phương lên bằng $25$? Có $5\cdot 5=25$, nhưng cũng có $(-5)\cdot(-5)=25$. Vậy $x=5$ hoặc $x=-5$.
\begin{warnbox}Rất nhiều bạn chỉ tìm được $x=5$ mà quên $x=-5$. Tích hai số âm là số dương, nên $(-5)^2=25$. Đề hỏi ``tất cả các số nguyên'' là một gợi ý: có thể có nhiều hơn một đáp số.\end{warnbox}
\verify{$x=-5$: $(-5)(-5)+(-5)(-5)=25+25=50$. Đúng.}
\ans{a) $14$;\ \ b) $x=4$;\ \ c) $x=5$ hoặc $x=-5$.}
}

% ---------------------------------------------------------------- Bài 34
\begin{probbox}{34}{\stars{1}}
\begin{enumerate}[label=\alph*),leftmargin=1.6em,itemsep=1pt,topsep=0pt]
\item Tổng của năm số lẻ liên tiếp bằng $155$. Tìm năm số đó.
\item Bạn Nam nói: ``Tớ tìm được bốn số lẻ liên tiếp có tổng bằng $150$.'' Nam nói đúng hay sai? Vì sao?
\end{enumerate}
\end{probbox}
\probsol{
\textbf{Phân tích đề.} Hai số lẻ liên tiếp hơn kém nhau $2$ (ví dụ $29$ và $31$). Câu a là bài toán tìm số; câu b yêu cầu \emph{giải thích}, không chỉ đoán.
\step{1}{Câu a: đặt ẩn là số ở giữa}
Gọi số ở giữa là $n$. Năm số là $n-4,\ n-2,\ n,\ n+2,\ n+4$.
\step{2}{Lập và giải phương trình}
$(n-4)+(n-2)+n+(n+2)+(n+4)=155$. Các số $-4$ và $+4$, $-2$ và $+2$ triệt tiêu nhau: $5n=155$, nên $n=31$.\\
Năm số là $27,\ 29,\ 31,\ 33,\ 35$.
\verify{$27+29+31+33+35=155$. Đúng, và đó đều là các số lẻ liên tiếp.}
\step{3}{Câu b: thử lập phương trình}
Gọi số lẻ bé nhất là $m$. Bốn số là $m,\ m+2,\ m+4,\ m+6$. Tổng: $4m+12=150$, suy ra $4m=138$, $m=138:4=34{,}5$.
Số $34{,}5$ không phải là số nguyên, càng không phải số lẻ. Vậy \emph{không có} bốn số lẻ liên tiếp nào có tổng $150$.
\begin{proofbox}Một cách giải thích khác bằng tính chẵn -- lẻ: tổng của hai số lẻ là số chẵn, nên tổng bốn số lẻ là số chẵn -- điều này chưa mâu thuẫn. Nhưng ta có thể đi xa hơn: $4m+12=4(m+3)$ luôn \emph{chia hết cho $4$}, trong khi $150$ không chia hết cho $4$ ($150=4\cdot 37+2$). Vậy tổng bốn số lẻ liên tiếp không bao giờ bằng $150$.\end{proofbox}
\begin{warnbox}Chỉ thử vài bộ số (như $35+37+39+41=152$) rồi kết luận là chưa đủ chặt chẽ. Ta cần một lí do đúng với \emph{mọi} trường hợp.\end{warnbox}
\ans{a) $27,\ 29,\ 31,\ 33,\ 35$;\ \ b) Nam nói sai, vì tổng bốn số lẻ liên tiếp luôn chia hết cho $4$ mà $150$ không chia hết cho $4$.}
}

% ---------------------------------------------------------------- Bài 35
\begin{probbox}{35}{\stars{1}}
Cứ $2$ thìa thức ăn thì đủ cho $12$ con cá ăn trong một ngày.
\begin{enumerate}[label=\alph*),leftmargin=1.6em,itemsep=1pt,topsep=2pt]
\item $5$ thìa thức ăn đủ cho bao nhiêu con cá ăn trong một ngày?
\item Bể cá nhà Hà có $18$ con cá. Một hộp thức ăn có $60$ thìa. Hộp thức ăn đó đủ cho cả bể ăn trong bao nhiêu ngày?
\end{enumerate}
\end{probbox}
\probsol{
\textbf{Phân tích đề.} Số cá được cho ăn trong một ngày tỉ lệ thuận với số thìa. Ta \emph{rút về đơn vị}: tính xem $1$ thìa đủ cho mấy con cá trong một ngày.
\begin{recallbox}Rút về đơn vị: tìm giá trị ứng với $1$ phần, rồi nhân lên (hoặc chia ra) theo số phần cần tìm.\end{recallbox}
\step{1}{Một thìa đủ cho bao nhiêu con cá một ngày?}
$12:2=6$ (con cá).
\step{2}{Câu a}
$5$ thìa: $5\cdot 6=30$ (con cá).
\step{3}{Câu b: mỗi ngày bể cần bao nhiêu thìa?}
Mỗi thìa đủ cho $6$ con trong một ngày, bể có $18$ con nên mỗi ngày cần $18:6=3$ (thìa).
\begin{figC}
\foreach \i in {0,1,2}{\draw[tmBlue!80!black,fill=tmBlue!12,rounded corners=2pt] (2.3*\i,0) rectangle (2.3*\i+2,0.6); \node[font=\footnotesize] at (2.3*\i+1,0.3){1 thìa: 6 con};}
\draw[decorate,decoration={brace,amplitude=5pt,mirror},Navy] (0,-0.1)--(6.6,-0.1) node[midway,below=6pt,font=\footnotesize]{$18$ con cá, mỗi ngày $3$ thìa};
\end{figC}
\step{4}{Số ngày}
Hộp có $60$ thìa, mỗi ngày dùng $3$ thìa: $60:3=20$ (ngày).
\begin{warnbox}Một lỗi hay gặp là lấy $60:18$ hoặc $60\cdot 6:18$ mà không hiểu ý nghĩa. Hãy luôn tự hỏi: ``Mỗi ngày dùng hết bao nhiêu thìa?'' rồi mới chia.\end{warnbox}
\verify{$20$ ngày $\times$ $3$ thìa $=60$ thìa, đúng bằng một hộp; và $3$ thìa $\times 6=18$ con mỗi ngày.}
\begin{parentbox}Có thể dùng đơn vị ``con-ngày'' cho dễ hình dung: mỗi thìa nuôi được $6$ ``con-ngày''; cả hộp nuôi được $60\cdot 6=360$ con-ngày; chia cho $18$ con được $20$ ngày.\end{parentbox}
\ans{a) $30$ con cá;\ \ b) $20$ ngày.}
}

% ---------------------------------------------------------------- Bài 36
\begin{probbox}{36}{\relax}
Mỗi lượt đi xe buýt, bạn Lan phải trả $7\,000$ đồng. Công ty xe buýt bán vé tháng giá $200\,000$ đồng, đi bao nhiêu lượt trong tháng cũng được.
\begin{enumerate}[label=\alph*),leftmargin=1.6em,itemsep=1pt,topsep=2pt]
\item Viết biểu thức biểu thị số tiền Lan phải trả nếu đi $n$ lượt trong tháng mà không mua vé tháng.
\item Trong một tháng, Lan phải đi ít nhất bao nhiêu lượt thì mua vé tháng sẽ rẻ hơn mua vé lượt?
\end{enumerate}
\end{probbox}
\probsol{
\textbf{Phân tích đề.} Mua vé lượt: số tiền phụ thuộc vào số lượt $n$. Mua vé tháng: luôn $200\,000$ đồng. Cần tìm số lượt nhỏ nhất để vé lượt \emph{đắt hơn} vé tháng.
\step{1}{Câu a}
Mỗi lượt $7\,000$ đồng, $n$ lượt: $7\,000\cdot n=7\,000n$ (đồng).
\step{2}{Câu b: so sánh hai cách}
Vé tháng rẻ hơn khi $7\,000n>200\,000$. Ta thử chia: $200\,000:7\,000=200:7\approx 28{,}57$.
\step{3}{Lập bảng thử những giá trị gần $28{,}57$}
\begin{center}
\begin{tabular}{|c|c|c|}
\hline
\rowcolor{tmBlue!10} Số lượt $n$ & Vé lượt $7\,000n$ (đồng) & So với vé tháng $200\,000$ đồng\\ \hline
$28$ & $196\,000$ & vé lượt rẻ hơn\\ \hline
$29$ & $203\,000$ & \textbf{vé tháng rẻ hơn}\\ \hline
\end{tabular}
\end{center}
Vậy cần đi ít nhất $29$ lượt.
\begin{warnbox}Số lượt phải là số tự nhiên, nên không trả lời ``$28{,}57$ lượt''. Cũng không ``làm tròn'' thành $28$: với $28$ lượt, vé lượt vẫn rẻ hơn. Ta phải lấy số tự nhiên nhỏ nhất \emph{lớn hơn} $28{,}57$, đó là $29$.\end{warnbox}
\begin{parentbox}Đây là bài toán ``đời thường'' rất tốt để trò chuyện với con: gói cước điện thoại, vé bơi theo tháng... đều có thể so sánh theo cách này.\end{parentbox}
\ans{a) $7\,000n$ đồng;\ \ b) ít nhất $29$ lượt.}
}

% ---------------------------------------------------------------- Bài 37
\begin{probbox}{37}{\relax}
Quan sát bảng sau (giá trị tăng đều khi $n$ tăng):
\begin{center}
\begin{tabular}{|c|c|c|c|c|c|}
\hline
\rowcolor{tmBlue!10} $n$ & $1$ & $2$ & $3$ & $4$ & $5$\\ \hline
Giá trị & $2$ & $7$ & $12$ & $17$ & $22$\\ \hline
\end{tabular}
\end{center}
\begin{enumerate}[label=\alph*),leftmargin=1.6em,itemsep=1pt,topsep=2pt]
\item Viết biểu thức theo $n$ cho hàng giá trị.
\item Tính giá trị ứng với $n=40$.
\item Với $n$ bằng bao nhiêu thì giá trị bằng $302$?
\item Có giá trị nào trong bảng (kéo dài mãi) bằng $2026$ không?
\end{enumerate}
\end{probbox}
\probsol{
\textbf{Phân tích đề.} Mỗi lần $n$ tăng $1$, giá trị tăng $5$. Đây là dãy cách đều với số hạng đầu $2$, khoảng cách $5$.
\begin{recallbox}Nếu $n$ tăng $1$ mà giá trị luôn tăng $d$ thì biểu thức có dạng $dn+b$; số hạng thứ $n$ của dãy cách đều là $a+(n-1)d$.\end{recallbox}
\step{1}{Câu a}
Biểu thức dạng $5n+b$. Thay $n=1$: $5+b=2$, nên $b=-3$. Biểu thức: $5n-3$.
(Hoặc: $2+5(n-1)=2+5n-5=5n-3$.)
\verify{$n=4$: $5\cdot 4-3=17$; \ $n=5$: $5\cdot 5-3=22$. Khớp bảng.}
\step{2}{Câu b}
$5\cdot 40-3=200-3=197$.
\step{3}{Câu c: giải phương trình}
$5n-3=302\ \Rightarrow\ 5n=305\ \Rightarrow\ n=61$.
\step{4}{Câu d}
$5n-3=2026\ \Rightarrow\ 5n=2029$. Nhưng $2029$ có chữ số tận cùng là $9$, không chia hết cho $5$, nên không có số tự nhiên $n$ nào thoả mãn.
\begin{proofbox}Mọi giá trị $5n-3$ đều có chữ số tận cùng là $2$ hoặc $7$ (vì $5n$ tận cùng là $0$ hoặc $5$). Số $2026$ tận cùng là $6$, nên không thể nằm trong bảng.\end{proofbox}
\ans{a) $5n-3$;\ \ b) $197$;\ \ c) $n=61$;\ \ d) không có.}
}

% ---------------------------------------------------------------- Bài 38
\begin{probbox}{38}{\stars{1}}
\begin{enumerate}[label=\alph*),leftmargin=1.6em,itemsep=1pt,topsep=0pt]
\item Tích $60\times 24\times 365$ biểu thị số đơn vị thời gian nào trong một năm không nhuận (có $365$ ngày)? Tính tích đó.
\item Một triệu giây bằng bao nhiêu ngày, giờ, phút, giây?
\end{enumerate}
\end{probbox}
\probsol{
\textbf{Phân tích đề.} Câu a: đọc tích như một chuỗi đổi đơn vị. Câu b: đổi ngược từ đơn vị nhỏ (giây) lên đơn vị lớn, bằng các phép chia có dư.
\begin{recallbox}$1$ phút $=60$ giây; $1$ giờ $=60$ phút $=3\,600$ giây; $1$ ngày $=24$ giờ $=24\cdot 3\,600=86\,400$ giây.\end{recallbox}
\step{1}{Câu a}
$365$ là số ngày trong năm; $24\times 365$ là số giờ trong năm; $60\times 24\times 365$ là số \textbf{phút} trong một năm không nhuận.
Tính: $24\times 365=8\,760$; \ $60\times 8\,760=525\,600$ (phút).
\step{2}{Câu b: đổi ra ngày}
$1\,000\,000:86\,400$ được thương $11$, dư $1\,000\,000-11\cdot 86\,400=1\,000\,000-950\,400=49\,600$ (giây).
\step{3}{Đổi phần dư ra giờ}
$49\,600:3\,600$ được thương $13$, dư $49\,600-13\cdot 3\,600=49\,600-46\,800=2\,800$ (giây).
\step{4}{Đổi phần dư ra phút và giây}
$2\,800:60$ được thương $46$, dư $2\,800-46\cdot 60=2\,800-2\,760=40$ (giây).
\begin{center}
\begin{tabular}{|l|r|r|}
\hline
\rowcolor{tmBlue!10} Phép chia & Thương & Số dư (giây)\\ \hline
$1\,000\,000:86\,400$ & $11$ ngày & $49\,600$\\ \hline
$49\,600:3\,600$ & $13$ giờ & $2\,800$\\ \hline
$2\,800:60$ & $46$ phút & $40$\\ \hline
\end{tabular}
\end{center}
\verify{$11\cdot 86\,400+13\cdot 3\,600+46\cdot 60+40=950\,400+46\,800+2\,760+40=1\,000\,000$.}
\begin{warnbox}Đừng viết kết quả dưới dạng thập phân như ``$11{,}57$ ngày'' rồi đọc thành ``$11$ ngày $57$ giờ'': phần thập phân của ngày không phải là số giờ. Hãy dùng phép chia có dư như trên.\end{warnbox}
\ans{a) Số phút trong một năm không nhuận, bằng $525\,600$;\ \ b) $11$ ngày $13$ giờ $46$ phút $40$ giây.}
}

% ---------------------------------------------------------------- Bài 39
\begin{probbox}{39}{\stars{1}}
Bốn số nguyên dương $a$, $b$, $c$, $d$ thoả mãn
\[\frac{3}{a-5}=\frac{3}{b+1}=\frac{3}{c-2}=\frac{3}{d+4}\qquad\text{và}\qquad a+b+c+d=50.\]
\begin{enumerate}[label=\alph*),leftmargin=1.6em,itemsep=1pt,topsep=2pt]
\item Sắp xếp $a$, $b$, $c$, $d$ theo thứ tự tăng dần.
\item Tìm giá trị của từng số.
\end{enumerate}
\end{probbox}
\probsol{
\textbf{Phân tích đề.} Bốn phân số cùng tử $3$ và bằng nhau. Dùng kĩ thuật ``đặt bằng $k$'' như Bài 12, rồi dùng thêm điều kiện về tổng để tìm $k$.
\step{1}{Các mẫu số bằng nhau}
\begin{recallbox}Nếu $\dfrac{m}{p}=\dfrac{m}{q}$ với $m\neq 0$ thì $p=q$.\end{recallbox}
$a-5=b+1=c-2=d+4$. Gọi giá trị chung là $k$ ($k\neq 0$ vì mẫu số phải khác $0$).
\step{2}{Biểu diễn theo $k$}
$a=k+5$,\quad $b=k-1$,\quad $c=k+2$,\quad $d=k-4$.
\step{3}{Câu a: so sánh}
$k-4<k-1<k+2<k+5$, nên $d<b<c<a$.
\step{4}{Câu b: dùng tổng để tìm $k$}
$a+b+c+d=(k+5)+(k-1)+(k+2)+(k-4)=4k+2$. Từ $4k+2=50$ suy ra $4k=48$, $k=12$.\\
Vậy $a=17$, \ $b=11$, \ $c=14$, \ $d=8$.
\verify{$a-5=12$, $b+1=12$, $c-2=12$, $d+4=12$: bốn phân số đều bằng $\dfrac{3}{12}$. Tổng $17+11+14+8=50$. Và $8<11<14<17$.}
\begin{warnbox}Khi cộng $(+5)+(-1)+(+2)+(-4)$ nhớ giữ đúng dấu: kết quả là $+2$, không phải $12$.\end{warnbox}
\ans{a) $d<b<c<a$;\ \ b) $a=17$, $b=11$, $c=14$, $d=8$.}
}

% ---------------------------------------------------------------- Bài 40
\begin{probbox}{40}{\stars{2}}
Có bốn túi gạo nặng khác nhau. Người ta cân từng cặp hai túi, đủ mọi cách ghép đôi, và ghi được sáu kết quả:
\[7 \text{ kg},\quad 8 \text{ kg},\quad 9 \text{ kg},\quad 10 \text{ kg},\quad 11 \text{ kg},\quad 12 \text{ kg}.\]
\begin{enumerate}[label=\alph*),leftmargin=1.6em,itemsep=1pt,topsep=2pt]
\item Tính tổng khối lượng của bốn túi.
\item Biết rằng khối lượng mỗi túi là một số nguyên ki-lô-gam. Tìm khối lượng của từng túi.
\end{enumerate}
\end{probbox}
\probsol{
\textbf{Phân tích đề.} Gọi khối lượng bốn túi từ nhẹ đến nặng là $a<b<c<d$ (kg). Bốn túi ghép đôi được $6$ cặp:
$a+b,\ a+c,\ a+d,\ b+c,\ b+d,\ c+d$. Đây là bài mở rộng của Bài 15 (ba quả bí).
\step{1}{Câu a: cộng cả sáu kết quả}
Mỗi túi có mặt trong đúng $3$ cặp (ghép với từng túi trong $3$ túi còn lại). Do đó
$3(a+b+c+d)=7+8+9+10+11+12=57$, \ nên $a+b+c+d=57:3=19$ (kg).
\begin{recallbox}Với ba vật, mỗi vật có mặt trong $2$ cặp nên tổng $=$ (tổng các cặp)$:2$. Với bốn vật, mỗi vật có mặt trong $3$ cặp nên tổng $=$ (tổng các cặp)$:3$.\end{recallbox}
\step{2}{Những cặp nào chắc chắn ứng với những kết quả nào?}
\begin{itemize}[leftmargin=1.4em,itemsep=1pt,topsep=1pt]
\item Cặp nhẹ nhất là hai túi nhẹ nhất: $a+b=7$.
\item Cặp nhẹ thứ hai là $a+c=8$: mọi cặp khác (trừ $a+b$) đều nặng hơn $a+c$, vì $a+d>a+c$, \ $b+c>a+c$, \ $b+d>a+c$, \ $c+d>a+c$.
\item Tương tự, cặp nặng nhất là $c+d=12$, cặp nặng thứ hai là $b+d=11$.
\item Hai kết quả ở giữa ($9$ và $10$) là $a+d$ và $b+c$, nhưng \emph{chưa biết} cái nào là $9$, cái nào là $10$.
\end{itemize}
\begin{proofbox}Lấy $a+c=8$ trừ $a+b=7$: $c-b=1$, tức $c=b+1$. Khi đó $b+c=b+(b+1)=2b+1$ là một \emph{số lẻ} (vì $b$ là số nguyên). Trong hai số $9$ và $10$, chỉ có $9$ là số lẻ. Vậy $b+c=9$ và $a+d=10$.\end{proofbox}
\step{3}{Tìm từng túi}
$2b+1=9\Rightarrow b=4$;\quad $c=b+1=5$;\quad $a=7-b=3$;\quad $d=12-c=7$.
\begin{barfigC}
\blk{0}{1.5}{0}{0.5}{tmGreen}{a=3}
\blk{1.5}{3.5}{0}{0.5}{tmOrange}{b=4}
\blk{3.5}{6}{0}{0.5}{tmPurple}{c=5}
\blk{6}{9.5}{0}{0.5}{tmBlue}{d=7}
\brcup{0}{9.5}{0.5}{tổng 19 kg}
\end{barfigC}
\verify{$3+4=7$, $3+5=8$, $3+7=10$, $4+5=9$, $4+7=11$, $5+7=12$: đúng sáu kết quả $7,8,9,10,11,12$. Tổng $3+4+5+7=19$.}
\begin{warnbox}Không được tự ý cho rằng kết quả thứ ba ($9$ kg) là $a+d$ hay $b+c$. Phải có lí do! Nếu không có điều kiện ``số nguyên'' thì trường hợp $b+c=10$ cho $a=2{,}5$; $b=4{,}5$; $c=5{,}5$; $d=6{,}5$ cũng thoả mãn.\end{warnbox}
\begin{parentbox}Có thể làm thí nghiệm bằng bốn mẩu giấy ghi số $3,4,5,7$: cho con ghép đôi tất cả các cách và cộng lại, con sẽ tự thấy mỗi số được dùng ba lần.\end{parentbox}
\ans{a) $19$ kg;\ \ b) $3$ kg, $4$ kg, $5$ kg, $7$ kg.}
}
